\begin{helper}
For every fixed \(I\), there is a finite history-cell family.  Its indices
\(i\) carry cells \(C_i\), representatives \(\omega_i\), and recorded
coordinate sets \(R_i\).  These cells cover exactly
\(\noderef{TwoBiteProjectionSizeEvent}(I,k,\ell_R,\ell_B)\), are pairwise
disjoint, contain their representatives, have the cylinder description
\[
\omega\in C_i
\Longleftrightarrow
\omega.2.2=\omega_i.2.2
\text{ and }
\forall e\in R_i,\ 
\noderef{TwoBiteEdgeCoordinateValue}(\omega,e)
\Longleftrightarrow
\noderef{TwoBiteEdgeCoordinateValue}(\omega_i,e),
\]
and are contained in the projection-size event.
\end{helper}
\begin{proof}
Fix deterministic orders of the finite base vertices, ordered base pairs, and
terminal coordinates.  For a sample \(\omega\), define its pre-terminal replay
as follows.  First record the injection.  From the injection compute projection
containment and \(F_1\).  Then scan the ordered \(F_2\)-candidate adjacencies.
For each same-color candidate that requires an edge value, insert the
corresponding tagged coordinate into the current recorded set and store its
\(\noderef{TwoBiteEdgeCoordinateValue}\).  Opposite-color candidates and loop
candidates require no graph coordinate.  After the \(F_2\)-scan, compute the
staged-revealed vertices, scan all incident coordinates from staged-revealed
vertices to projection-contained endpoints, and finally scan all projected
pairs coming from the \(X_x(I)\)-sets of staged-revealed vertices.  Each
queried terminal coordinate is inserted into the recorded set before its truth
value is used.

The transcript is the injection, the final recorded set, and the truth table
on that final recorded set.  There are only finitely many possible
transcripts: the injection has finite domain and codomain, the final recorded
set is a subset of the finite terminal-coordinate universe, and a truth table
on a finite set has finitely many values.  Index the transcripts which are
realized by samples in
\(\noderef{TwoBiteProjectionSizeEvent}(I,k,\ell_R,\ell_B)\), and choose one
representative \(\omega_i\) for each realized transcript.  Let \(R_i\) be the
final recorded set in the transcript of \(\omega_i\).

The key point is that the cylinder condition in the statement is not weaker
than equality of transcripts.  We prove this by induction through the replay
run.  Suppose \(\omega\) has the same injection as \(\omega_i\) and agrees
with \(\omega_i\) on every coordinate in \(R_i\).  Projection containment and
\(F_1\) depend only on the injection, so they agree at the start of the replay.
During the \(F_2\)-scan, the two runs therefore consider the same ordered
candidates.  If a candidate has opposite colors or is a loop, both runs make
the same deterministic decision without reading a coordinate.  If it is a
same-color candidate, the coordinate read in the representative run is
inserted into \(R_i\) before its value is used.  This is the same class of
\(F_2\)-scan coordinates isolated by
\(\noderef{FixedSetPreTerminalF2QueriesRecorded}\).  Agreement on \(R_i\)
gives the same edge value, so the candidate contributes the same way to the
\(F_2\)-count in both runs.  Hence \(F_2\), and then the staged-revealed set,
agree.

The incident-coordinate pass is now the same in both runs, because it is
indexed by the common projection data and the common staged-revealed set.
Every incident coordinate read by the representative replay lies in \(R_i\),
so the truth-table agreement gives the same incident edge values.  For the
final \(X_x(I)\)-pair pass, fix a staged-revealed base vertex \(x\) and
\(v\in I\).  Membership of \(v\) in \(X_x(I)\) is membership in the lifted
neighborhood of \(x\).  In the red case this is the red graph edge from the
red endpoint of \(x\) to the red projection of \(v\), and in the blue case it
is the analogous blue edge.  The projection of \(v\) is the same in the two
samples because the injections agree, and this incident edge was read in the
previous pass because \(x\) is staged-revealed and the projected endpoint is
projection-contained.  Thus it is in \(R_i\), so its truth value agrees.
Therefore all \(X_x(I)\)-sets agree, all projected pair coordinates considered
in the final pass agree, and every final queried coordinate again receives the
same truth value.  The two replays consequently have the same final transcript.

Thus the set of samples with the same transcript as \(\omega_i\) is exactly
the cylinder described in the statement: common injection and common truth
table on \(R_i\).  Define \(C_i\) to be this cylinder.  The cells cover the
projection-size event because every sample in the event has a realized replay
transcript and hence belongs to the cell of its representative.  They are
pairwise disjoint because the deterministic replay produces a unique
transcript, and the preceding paragraph identifies each cylinder with one
transcript class.  Each representative lies in its own cell by reflexivity.
Finally, \(\noderef{TwoBiteProjectionSizeEvent}\) depends only on the two
projection images of \(I\), which are determined by the injection; therefore
every sample in the cell of a representative chosen from the event also lies
in the projection-size event.
\end{proof}
